% ═══════════════════════════════════════════════════════════════
% AB-01a: Grundgesetz der Mechanik (F = m · a)
% Physik Klasse 10 MR
% ═══════════════════════════════════════════════════════════════

\documentclass[11pt, a4paper]{article}

\usepackage[utf8]{inputenc}
\usepackage[T1]{fontenc}
\usepackage[ngerman]{babel}

\usepackage{helvet}
\renewcommand{\familydefault}{\sfdefault}

\usepackage{microtype}
\usepackage[margin=1.5cm, top=2cm, bottom=1.5cm]{geometry}
\usepackage{enumitem}
\usepackage{tabularx}
\usepackage{booktabs}
\usepackage{array}
\usepackage{tikz}
\usetikzlibrary{calc}

\usepackage{amsmath, amssymb}
\usepackage{siunitx}
\sisetup{locale = DE, per-mode = symbol}

\usepackage{fancyhdr}
\usepackage{lastpage}
\usepackage{needspace}

\usepackage{xcolor}
\usepackage{tcolorbox}
\tcbuselibrary{skins}

% ═══════════════════════════════════════════════════════════════
% FARBEN
% ═══════════════════════════════════════════════════════════════

\definecolor{physik}{HTML}{2c5aa0}
\definecolor{hellgrau}{HTML}{F5F5F5}
\definecolor{mittelgrau}{HTML}{888888}
\definecolor{rahmen}{HTML}{CCCCCC}

\colorlet{fachfarbe}{physik}

% ═══════════════════════════════════════════════════════════════
% VARIABLEN
% ═══════════════════════════════════════════════════════════════

\newcommand{\fach}{Physik}
\newcommand{\thema}{Grundgesetz der Mechanik}
\newcommand{\klasse}{10 MR}
\newcommand{\stunde}{01a}
\newcommand{\maxpunkte}{28}

% ═══════════════════════════════════════════════════════════════
% KOPF- UND FUSSZEILE
% ═══════════════════════════════════════════════════════════════

\pagestyle{fancy}
\fancyhf{}
\renewcommand{\headrulewidth}{0.4pt}
\renewcommand{\footrulewidth}{0pt}
\lhead{\small \fach{} Klasse \klasse{} | \thema}
\rhead{\small Name: \underline{\hspace{3.5cm}}}
\cfoot{\small\textcolor{mittelgrau}{Seite \thepage{} von \pageref{LastPage}}}

% ═══════════════════════════════════════════════════════════════
% BOXEN
% ═══════════════════════════════════════════════════════════════

\newtcolorbox{merkkasten}[1][]{
    colback=hellgrau,
    colframe=hellgrau,
    coltitle=black,                          % WICHTIG: Titel-Text schwarz (nicht Standard-weiß!)
    colbacktitle=hellgrau,                   % Titel-Hintergrund wie Box
    boxrule=0pt,
    arc=0pt,
    left=8pt, right=8pt, top=8pt, bottom=6pt,
    borderline north={2pt}{0pt}{fachfarbe},
    before skip=6pt, after skip=6pt,
    fonttitle=\bfseries,                     % Titel fett
    #1
}

\newtcolorbox{formelkasten}{
    colback=white,
    colframe=fachfarbe,
    boxrule=1.5pt,
    arc=0pt,
    left=8pt, right=8pt, top=6pt, bottom=6pt,
    before skip=4pt, after skip=4pt
}

% ═══════════════════════════════════════════════════════════════
% BEFEHLE
% ═══════════════════════════════════════════════════════════════

\newcommand{\luecke}[1]{\underline{\hspace{#1}}}
\newcommand{\punktebox}[1]{\hfill\small\textcolor{mittelgrau}{[\_\_/#1]}}

\newcommand{\aufgabe}[1]{%
    \needspace{4\baselineskip}%
    \vspace{10pt}%
    \noindent\textbf{\textcolor{fachfarbe}{Aufgabe #1}}%
    \vspace{4pt}%
    \nopagebreak[4]%
}

\newcommand{\operator}[1]{\textbf{\textcolor{fachfarbe}{#1}}}

\newlist{teil}{enumerate}{2}
\setlist[teil,1]{label=\alph*), leftmargin=1.8em, itemsep=3pt, parsep=0pt, topsep=3pt}

\newcommand{\antwortzeilen}[1]{%
    \par\vspace{4pt}%
    \foreach \n in {1,...,#1}{%
        \noindent\rule{\linewidth}{0.3pt}\vspace{6pt}%
    }%
}

\setlength{\parindent}{0pt}
\setlength{\parskip}{4pt}

% ═══════════════════════════════════════════════════════════════
% DOKUMENT
% ═══════════════════════════════════════════════════════════════

\begin{document}

% ═══════════════════════════════════════════════════════════════
% HEADER
% ═══════════════════════════════════════════════════════════════

\begin{center}
    {\Large\bfseries\textcolor{fachfarbe}{\thema: $F = m \cdot a$}}\\[0.2em]
    {\small Arbeitsblatt | \fach{} Klasse \klasse{} | Stunde \stunde}
\end{center}
\vspace{-0.3em}
\hrule
\vspace{0.6em}

% ═══════════════════════════════════════════════════════════════
% KASTEN 1: RESULTIERENDE KRAFT
% ═══════════════════════════════════════════════════════════════

\begin{merkkasten}[title={\textbf{Kasten 1: Resultierende Kraft}}]
Die \textbf{resultierende Kraft} ist die \luecke{3cm} aller auf einen Körper wirkenden Kräfte.

\vspace{0.3em}
Kräfte in dieselbe Richtung werden \luecke{2.5cm}.

Kräfte in entgegengesetzte Richtung werden \luecke{2.5cm}.
\end{merkkasten}

% ═══════════════════════════════════════════════════════════════
% KASTEN 2a: ERKENNTNIS KRAFT
% ═══════════════════════════════════════════════════════════════

\begin{merkkasten}[title={\textbf{Kasten 2a: Erkenntnis – Kraft und Beschleunigung}}]
Zwei Kisten mit \textbf{gleicher Masse} werden mit \textbf{unterschiedlicher Kraft} geschoben.

\vspace{0.3em}
\textbf{Beobachtung:} Die Kiste mit der größeren Kraft beschleunigt \luecke{2.5cm}.

\vspace{0.3em}
\textbf{Erkenntnis:} Je größer die \textbf{Kraft}, desto \luecke{2.5cm} die Beschleunigung.
\end{merkkasten}

% ═══════════════════════════════════════════════════════════════
% KASTEN 2b: ERKENNTNIS MASSE
% ═══════════════════════════════════════════════════════════════

\begin{merkkasten}[title={\textbf{Kasten 2b: Erkenntnis – Masse und Beschleunigung}}]
Zwei Einkaufswagen mit \textbf{unterschiedlicher Masse} werden mit \textbf{gleicher Kraft} geschoben.

\vspace{0.3em}
\textbf{Beobachtung:} Der schwere Wagen beschleunigt \luecke{2.5cm}.

\vspace{0.3em}
\textbf{Erkenntnis:} Je größer die \textbf{Masse}, desto \luecke{2.5cm} die Beschleunigung.
\end{merkkasten}

% ═══════════════════════════════════════════════════════════════
% KASTEN 2c: VON DER BEOBACHTUNG ZUR FORMEL
% ═══════════════════════════════════════════════════════════════

\begin{merkkasten}[title={\textbf{Kasten 2c: Von der Beobachtung zur Formel}}]
\textbf{Aus Kasten 2a:} Mehr Kraft $\rightarrow$ mehr Beschleunigung

\hspace{2em}Mathematisch: $a \sim F$ \quad {\small(a ist \luecke{3cm} zu F)}

\vspace{0.4em}
\textbf{Aus Kasten 2b:} Mehr Masse $\rightarrow$ weniger Beschleunigung

\hspace{2em}Mathematisch: $a \sim \dfrac{1}{m}$ \quad {\small(a ist \luecke{3cm} proportional zu m)}

\vspace{0.5em}
\hrule
\vspace{0.5em}

\textbf{Beide zusammen:} \quad $a = \dfrac{F}{m}$ \quad $\Rightarrow$ \quad umgestellt: \quad $F = $ \luecke{2cm}
\end{merkkasten}

% ═══════════════════════════════════════════════════════════════
% KASTEN 3: GRUNDGESETZ DER MECHANIK
% ═══════════════════════════════════════════════════════════════

\begin{formelkasten}
\begin{center}
\textbf{Kasten 3: Das Grundgesetz der Mechanik (Newton)}

\vspace{0.5em}
{\Large $F = m \cdot a$}

\vspace{0.5em}
\begin{tabular}{lll}
\textbf{F} = & \luecke{3cm} & Einheit: \luecke{2cm} \\[0.5em]
\textbf{m} = & \luecke{3cm} & Einheit: \luecke{2cm} \\[0.5em]
\textbf{a} = & \luecke{3cm} & Einheit: \luecke{2cm} \\
\end{tabular}
\end{center}
\end{formelkasten}

% ═══════════════════════════════════════════════════════════════
% KASTEN 4: FORMEL-DREIECK
% ═══════════════════════════════════════════════════════════════

\begin{merkkasten}[title={\textbf{Kasten 4: Das Formel-Dreieck}}]
\begin{minipage}{0.45\linewidth}
\begin{center}
\begin{tikzpicture}[scale=0.9]
    % Dreieck
    \draw[fachfarbe, line width=1.5pt] (0,0) -- (3,0) -- (1.5,2.2) -- cycle;
    % Horizontale Linie
    \draw[fachfarbe, line width=1pt] (0.45,1.1) -- (2.55,1.1);
    % Vertikale Linie
    \draw[fachfarbe, line width=1pt] (1.5,0) -- (1.5,1.1);
    % Beschriftung
    \node at (1.5,1.6) {\Large\textbf{F}};
    \node at (0.75,0.5) {\Large\textbf{m}};
    \node at (2.25,0.5) {\Large\textbf{a}};
\end{tikzpicture}
\end{center}
\end{minipage}%
\hfill
\begin{minipage}{0.5\linewidth}
\textbf{Die drei Formeln:}

\vspace{0.3em}
$F = $ \luecke{2cm}

\vspace{0.3em}
$m = $ \luecke{2cm}

\vspace{0.3em}
$a = $ \luecke{2cm}
\end{minipage}
\end{merkkasten}

% ═══════════════════════════════════════════════════════════════
% AUFGABEN
% ═══════════════════════════════════════════════════════════════

\aufgabe{1: Kraft berechnen (F = m $\cdot$ a)} \punktebox{6}

\operator{Berechne} die Kraft in den folgenden Situationen.

\begin{teil}
    \item Ein Skateboard ($m = 5\,$kg) wird mit $a = 2\,\text{m/s}^2$ beschleunigt.

    \vspace{0.5em}
    $F = $ \luecke{6cm} $= $ \luecke{2cm}
    \vspace{0.8em}

    \item Ein Fahrrad ($m = 20\,$kg) beschleunigt mit $a = 5\,\text{m/s}^2$.

    \vspace{0.5em}
    $F = $ \luecke{6cm} $= $ \luecke{2cm}
    \vspace{0.8em}

    \item Ein Auto ($m = 1000\,$kg) beschleunigt mit $a = 4\,\text{m/s}^2$.

    \vspace{0.5em}
    $F = $ \luecke{6cm} $= $ \luecke{2cm}
    \vspace{0.5em}
\end{teil}

% ---

\aufgabe{2: Masse berechnen (m = F / a)} \punktebox{6}

\operator{Berechne} die Masse.

\begin{teil}
    \item $F = 50\,$N, $a = 5\,\text{m/s}^2$. Wie schwer ist der Koffer?

    \vspace{0.3em}
    $m = $ \luecke{6cm} $= $ \luecke{2cm}

    \item $F = 200\,$N, $a = 10\,\text{m/s}^2$

    \vspace{0.3em}
    $m = $ \luecke{6cm} $= $ \luecke{2cm}

    \item $F = 400\,$N, $a = 8\,\text{m/s}^2$

    \vspace{0.3em}
    $m = $ \luecke{6cm} $= $ \luecke{2cm}
\end{teil}

% ---

\aufgabe{3: Beschleunigung berechnen (a = F / m)} \punktebox{6}

\operator{Berechne} die Beschleunigung.

\begin{teil}
    \item Ein Ball ($m = 2\,$kg) wird mit $F = 10\,$N angestoßen.

    \vspace{0.3em}
    $a = $ \luecke{6cm} $= $ \luecke{2cm}

    \item $F = 100\,$N, $m = 20\,$kg

    \vspace{0.3em}
    $a = $ \luecke{6cm} $= $ \luecke{2cm}

    \item Ein Motorrad ($m = 200\,$kg) wird mit $F = 1000\,$N beschleunigt.

    \vspace{0.3em}
    $a = $ \luecke{6cm} $= $ \luecke{2cm}
\end{teil}

% ---

\aufgabe{4: Resultierende Kraft} \punktebox{6}

Ein Auto ($m = 500\,$kg) wird von zwei Kräften beeinflusst:
\begin{itemize}[leftmargin=2em, itemsep=2pt]
    \item Motor schiebt nach vorne: $F_{\text{Motor}} = 3000\,$N
    \item Reibung bremst nach hinten: $F_{\text{Reib}} = 500\,$N
\end{itemize}

\begin{teil}
    \item \operator{Berechne} die resultierende Kraft.

    \vspace{0.3em}
    $F_{\text{res}} = F_{\text{Motor}} - F_{\text{Reib}} = $ \luecke{4cm} $= $ \luecke{2cm}

    \item \operator{Berechne} die Beschleunigung des Autos.

    \vspace{0.3em}
    $a = \dfrac{F_{\text{res}}}{m} = $ \luecke{4cm} $= $ \luecke{2cm}

    \item \operator{Erkläre}, warum das Auto beschleunigt (nicht bremst).

    \antwortzeilen{2}
\end{teil}

% ---

\aufgabe{5: Transfer} \punktebox{4}

Beim Tauziehen ziehen 4 Personen links mit je $200\,$N und 3 Personen rechts mit je $300\,$N.

\begin{teil}
    \item \operator{Berechne} die resultierende Kraft und ihre Richtung.

    \vspace{0.3em}
    $F_{\text{links}} = $ \luecke{2cm}, \quad $F_{\text{rechts}} = $ \luecke{2cm}

    \vspace{0.3em}
    $F_{\text{res}} = $ \luecke{4cm} Richtung: \luecke{2cm}

    \item Wenn das Seil mit beiden Teams $m = 100\,$kg wiegt: Wie groß ist die Beschleunigung?

    \vspace{0.3em}
    $a = $ \luecke{6cm} $= $ \luecke{2cm}
\end{teil}

% ═══════════════════════════════════════════════════════════════
% PUNKTE
% ═══════════════════════════════════════════════════════════════

\vfill
\hrule
\vspace{0.4em}

\hfill\textbf{Punkte:} \luecke{1.2cm} / \maxpunkte

\end{document}
